% ============================================================================
% Fichier  : front/01_avant_propos.tex
% Titre    : Avant-propos
% Auteur   : MINKA MI NGUIDJOÏ Thierry Emmanuel
% Version  : 1.0 -- novembre 2026
% ============================================================================

\chapter*{Avant-propos}
\addcontentsline{toc}{chapter}{Avant-propos}
\markboth{Avant-propos}{Avant-propos}

\epigraph{« Ce qui distingue le testeur de l'attaquant,\\
ce n'est pas l'outil~: c'est l'autorisation. »}{Principe du cours}

Ce cours s'adresse aux élèves ingénieurs de quatrième année du Département
de Génie Informatique de l'École Nationale Supérieure Polytechnique de
Yaoundé. Il suppose acquis le cours de sécurité informatique (SEC3266), en
particulier les réseaux, les identités, les systèmes et les applications, et
s'articule avec l'audit des systèmes d'information (INF3226) et l'analyse
des risques (GES4062). Il est la seconde partie de l'enseignement du hacking
éthique~: il approfondit la démarche du test d'intrusion, des règles
d'engagement au rapport.

Il repose sur une conviction~: on ne défend bien que ce que l'on sait
attaquer. Mais il en tire une exigence~: chaque technique est enseignée avec
son cadre légal, sa détection et sa parade, et uniquement pour des systèmes
que l'on est autorisé à tester. Le cours ne fournit pas de code d'attaque
prêt à l'emploi~; il apprend à raisonner comme un attaquant pour conseiller
un défenseur.

\bigskip

\textbf{Ce que ce cours n'est pas.} Il n'est pas une préparation à une
certification, ni un catalogue d'outils, ni une invitation à tester sans
mandat. Un test non autorisé est une infraction au Cameroun, même animé par
la curiosité~: la loi vise expressément l'accès obtenu « par défi
intellectuel ».

\bigskip

\textbf{Ce que ce cours est.} Un traité de vingt-six chapitres, en six
parties précédées d'un prologue. La Partie~I pose le cadre~: définitions et
métiers, droit camerounais, règles d'engagement, méthodologies. La
Partie~II apprend à reconnaître~: renseignement en sources ouvertes,
balayage, analyse des vulnérabilités. La Partie~III traite des systèmes et
des réseaux, dont l'annuaire Active Directory. La Partie~IV porte sur les
applications web, les API, le mobile et le nuage. La Partie~V va au-delà des
systèmes~: sans fil, ingénierie sociale, équipe rouge, détection. La
Partie~VI apprend à restituer~: rapport, contre-test, divulgation,
laboratoire, mesure.

\bigskip

\textbf{Un fil rouge.} Après la crise de mars (Opération GAROUA, cours
SEC3266 et GES4062), LiZbiZ commande un test d'intrusion de son système
(Opération LOGONE). La BANQUE SANAGA conduit un exercice d'équipe rouge
(Opération MBAM), et l'ANPP fait tester son portail et son API (Opération
MEFOU). Les noms, les dates, les adresses et les montants sont fictifs~; les
mécanismes sont exacts, et les textes de loi sont cités d'après leur version
officielle.

\bigskip

\textbf{Comment travailler.} Lisez chaque chapitre avec sa fiche
d'auto-évaluation~; faites les exercices avant de lire leurs corrigés~;
pratiquez dans les labs de NDAP SUKULU, qui travaillent sur des pièces et
jamais sur des systèmes réels. Pour aller plus loin, construisez votre
propre laboratoire isolé (Chapitre~\ref{chap:labo}).

\bigskip

\begin{flushright}
    \IfFileExists{../../Logos/signature_avant_propos.png}{\includegraphics[height=1.6cm]{../../Logos/signature_avant_propos.png}}{\textit{MINKA MI NGUIDJOÏ Thierry Emmanuel}}\\
    \textit{Yaoundé, novembre 2026}
\end{flushright}
